%Hier den Inhalt mit \input einfügen und die folgenden Zeilen entfernen

%\input{nomenclature.tex}

\input{01_Wiederholung}
\newpage

\input{02_Frequenzgang}
\newpage

\input{03_LogDarstellung}
\newpage

\input{04_Ortskurven}
\newpage

\input{05_Resonanzkreise}

%%%%%%%%%%%%%%%%%%%%%%%%%%%%%%%%%
% NOTIZEN ZUM INHALT DES MODULS %
%%%%%%%%%%%%%%%%%%%%%%%%%%%%%%%%%

% Allgemeiner Aufbau:
% -------------------
% Kapitel 01 - Modulüberblick und Wiederholung
% Kapitel 02 - Frequenzgang, Amplitudengang, Phasengang
% Kapitel 03 - Logarithmische Darstellung Frequenzgang
% Kapitel 04 - Ortskurven
% Kapitel 05 - Resonanzkreise, Schwingungen
% Kapitel 06 - Netzwerkberechnung (discarded)


%%%%%%%%%%%%%%%%%%%%%%%%%%%%%%%%%%%%%%%%%%%%%%%%%%%%%%%%%%%%%%%%%%%%%%%%%%%%%%%%%%%
%         D U M P   Z O N E   ---    T E S T S   A N D   T R Y O U T S            %
%%%%%%%%%%%%%%%%%%%%%%%%%%%%%%%%%%%%%%%%%%%%%%%%%%%%%%%%%%%%%%%%%%%%%%%%%%%%%%%%%%%
\s{% Skript only
\section{Nebenrechnungen}
\begin{frame}{\secname}

    % Nebenrechnung 1
    Frequenzgang eines RC-Tiefpass 1. Ordnung. [Nebenrechnung zu Gleichung \eqref{eq:tiefpass:frequenzgang}]
    \begin{equation}
        \begin{aligned}
            \underline{F}(\mathrm{j}\omega)
                & = \frac{\underline{U}_2}{\underline{U}_1}
                    %&&\left|\ \right.
                    &&\left|\ =\frac{Z_C}{\left(Z_R + Z_C\right)} \frac{\cdot\redcancel{\underline{I}}}{\cdot\redcancel{\underline{I}}}\right.
                    &&\ \text{Komplexe Spannungsteilerregel}\\
                &= \frac{\frac{1}{\mathrm{j}\omega C}}{R + \frac{1}{\mathrm{j}\omega C}}
                    &&\left|\ \cdot \frac{\mathrm{j}\omega C}{\mathrm{j}\omega C}\right.
                    &&\ \text{Rationalisieren}\\
                &= \frac{1}{\mathrm{j}\omega CR + 1}
                    &&\left|\ \vphantom{\frac{0}{0}}\right.
                    &&\ \text{Umformen}\\
                &= \frac{1}{1 + \mathrm{j}\omega CR}
                    &&\left|\ \cdot \frac{1 - \mathrm{j}\omega CR}{1 - \mathrm{j}\omega CR}\right.
                    &&\ \text{Optional für kartesische Form}\\
                &= \frac{1 - \mathrm{j}\omega CR}{1 + (\omega CR)^2}
                % Alternativ \Big| oder \Bigg| für größere vertikale Striche (absolute Größe)
            \label{eq:tiefpass:frequenzgang:nebenrechnung}
        \end{aligned}
    \end{equation}
\end{frame}
}
%%%%%%%%%%%%%%%%%%%%%%%%%%%%%%%%%%%%%%%%%%%%%%%%%%%%%%%%%%%%%%%%%%%%%%%%%%%%%%%%%%%%%%%%%%%%%%%%%%%
\CMT{
\section{NOTES}
\begin{frame}{\secname}

    Notes:
    \begin{itemize}
        \item https://www.falstad.com/afilter/ Circuit Simulator (Amplitudengang, Auswahl typischer Filter)
        \item[] e.g. frame, itemize, column, tikzpicture, figure with subfigure, minipage, etc.
        \item[] Weiterführende Quellen
    \end{itemize}
\end{frame}
}
%%%%%%%%%%%%%%%%%%%%%%%%%%%%%%%%%%%%%%%%%%%%%%%%%%%%%%%%%%%%%%%%%%%%%%%%%%%%%%%%%%%%%%%%%%%%%%%%%%%

% Beispiel für Bandsperre als Parallelschaltung von Tiefpass und Hochpass (Innerhalb des \CMT{} Kommentars)
% https://www.falstad.com/afilter/
\CMT{
$ 3 0.000005 5 50 5 50
% 0 5290848.156670315
w 264 128 248 128 0
w 200 56 200 128 0
O 264 56 408 56 0
170 200 56 168 56 2 20 4000 5 0.1
c 200 128 248 128 0 0.000001 0
c 264 56 264 128 0 0.000001 0
r 200 56 264 56 0 100
r 264 128 264 216 0 1
g 264 216 264 248 0
}
% RLC Bandpass:
% https://www.falstad.com/circuit/e-bandpass.html
% In Moodle als iframe einbindbar (via weburl oder selbst gehosted)

%%%%%%%%%%%%%%%%%%%%%%%%%%%%%%%%%%%%%%%%%%%%%%%%%%%%%%%%%%%%%%%%%%%%%%%%%%%%%%%%%%%%%%%%%%%%%%%%%%%
%        L I T E R A T U R   U N D   Q U E L L E N A N G A B E N                                  %
%%%%%%%%%%%%%%%%%%%%%%%%%%%%%%%%%%%%%%%%%%%%%%%%%%%%%%%%%%%%%%%%%%%%%%%%%%%%%%%%%%%%%%%%%%%%%%%%%%%

% \section{Literatur und Quellen}
% \begin{frame}
% \s{\cleardoublepage}
% \input{literatur.tex}
% \end{frame}